% ============================================================
% Chapter 2: State of the Art
%   Section: News Discovery, Selection and Recommendation
% Included from chapters/02-state-of-the-art/state-of-the-art.tex
% Context: docs/arquitectura-tecnica.md section 3 (the recommendation
% system of the wider project, out of scope here).
% ============================================================

\section{News Discovery, Selection and Recommendation}
\label{sec:sota-selection}

Before a claim can be checked, someone has to decide which articles are
worth reading. Two traditions address that decision.

\paragraph{Editorial selection.} In a newsroom, an editor reads what
the sources published and chooses what to run; aggregators of
constructive news do the same by hand. The decision is a judgement of
relevance and of fit with the publication's line, and it is the step
this work tries to support rather than personalise: the system reads
the sources, screens out what the publication never covers, ranks what
is left, and leaves the choice to a person or to a language model that
reads titles and summaries. Language models have been reported to match
or exceed crowd workers on several text-annotation tasks
\cite{gilardi2023chatgpt}, which makes them plausible stand-ins for a
reader's first pass; whether they agree with an editor on \emph{this}
judgement is part of what RQ3 measures.

\paragraph{Recommendation.} News recommender systems choose, per
reader, from a large and fast-changing pool. Their survey by Karimi,
Jannach and Jugovac lists what makes news a hard domain --- short item
lifetimes, the importance of recency, popularity effects, and the
risk of narrowing what each reader sees \cite{karimi2018news} --- the
last of which Pariser called the filter bubble
\cite{pariser2011filter}. Industrial systems combine content
similarity with graph structure: Twitter's SimClusters represents users
and items by the communities they belong to
\cite{satuluri2020simclusters}, and TwHIN learns embeddings over a
heterogeneous graph of users, items and topics
\cite{elkishky2022twhin}.

The wider project behind this thesis plans a recommender of the second
kind over the graph that every analysis writes, with the verification
verdict as one of its signals (docs/arquitectura-tecnica.md,
section~3). That recommender is out of scope here
(Section~\ref{sec:scope}); this thesis builds the editorial selection
that comes before it and the graph that would feed it.
